\documentclass[aspectratio=169,xcolor=dvipsnames]{beamer}
\usetheme{SimpleDarkBlue}
\usepackage{hyperref}
\usepackage{graphicx} % Allows including images
\usepackage{booktabs} % Allows the use of \toprule, \midrule and \bottomrule in tables
\usepackage[spanish]{babel}
\usepackage[utf8]{inputenc}
\usepackage{array}
\usepackage{tabularx}
\usepackage{caption}
\usepackage{adjustbox}
\usepackage{amssymb}
\usepackage{tcolorbox}
\usepackage{amsmath}
\usepackage[T1]{fontenc}

\graphicspath{{./graficas}} 

\title{Estadística aplicada}
\subtitle{Conceptos previos: conjuntos, funciones, cálculo y probabilidad}

\author{Práxedis Jiménez Ruvalcaba}

\institute
{

}
\date{\today}

\tcbuselibrary{theorems}

% Custom colored theorem box
\newtcbtheorem{mytheo}{Teorema}%
{colback=orange!5!white,colframe=orange!75!black,fonttitle=\bfseries}{th}

\begin{document}

\begin{frame}[c]
    \begin{beamercolorbox}[sep=8pt,center,rounded=true,shadow=true]{title}
        \usebeamerfont{title}\inserttitle\par
        \ifx\insertsubtitle\@empty
        \else
          \vskip0.25em
          {\usebeamerfont{subtitle}\insertsubtitle\par}
        \fi 
    \end{beamercolorbox}

    \vspace{1cm}
    \begin{center}
        \usebeamerfont{author}\insertauthor\par
        \vspace{0.5cm}
        \usebeamerfont{institute}\insertinstitute\par
        \usebeamerfont{date}\insertdate
    \end{center}
    
\end{frame}

\section{Diferencias entre tipo}

\begin{frame}{Tipos de datos}
    A nivel práctico de manejo de R, podemos considerar que R solo cuenta con:
    \begin{itemize}
        \item double (númerico)
        \item character 
        \item bool
    \end{itemize}
    R si cuenta con datos tipo integer, string, entre otros. Pero para efectos prácticos, 
    solo se operan los anteriores
\end{frame}

\begin{frame}{Tipos de variables}
        \begin{center}
        \begin{tabular}{|c|p{11cm}|}
            \hline
            \textbf{Tipo de dato} & \textbf{Característica} \\ \hline
            Vector & Guarda los datos de un mismo tipo. Se le aplican la mayoría de funciones a nivel operativo. \\
            Factor & Para categorías y contruir datos categóricos \\
            DataFrame & Conjunto de datos, la manera estándar de operar la 'data' \\ 
            Table & Mayormente utilizada para frecuencias, dependiente de las categorías \\
            Lists & Guardar datos de varios tipos distintos \\
            Matrix & Un herramienta de manejo de datos a nivel tensorial \\
            \hline
        \end{tabular}
    \end{center}
\end{frame}

\section{Comandos importantes}

\begin{frame}{Comandos que nunca vienen mal}
    \begin{center}
        \begin{tabular}{|c|p{11cm}|}
            \hline
            \textbf{Comando} & \textbf{Funcionalidad} \\ \hline
            SessionInfo() & Brinda los detalles completo de lo que se está utilizando \\
            ?\textit{función} & Muestra la documentación interna sobre \textit{función} \\
            ??\textit{función} & Muestra todas las funciones con nombres similares o dónde aparezca menionada \textit{función} \\
            typeof() & Descripe el tipo de dato o variable \\

            \hline
        \end{tabular}
    \end{center}

\end{frame}

\begin{frame}{Comandos para dataframes}
    si se tiene un dataframe '\textit{df}', entonces:
    \begin{center}
        \begin{tabular}{|c|p{9cm}|}
            \hline
            \textbf{Comando} & \textbf{Funcionalidad} \\ \hline
            rbind(df,otro\_df) & A df le concatena otro\_df \\
            duplicated(df) & Obtiene un vector booleano de qué renglones son duplicados \\
            order(df \texttt{\char36} columna) & Obtiene un vector con el orden de las filas según una columna \\ 
            \hline
        \end{tabular}
    \end{center}
\end{frame}


\end{document}